\documentclass[11pt]{article}
\usepackage[margin=1in]{geometry}
\usepackage{amsmath}
\usepackage{amssymb}
\usepackage{parskip}
\usepackage{xcolor}
\colorlet{green}{green!55!black}
\usepackage[hidelinks]{hyperref}

\title{Non-stationary conditional extremes for BC streamflow:\\ my current understanding of the model}
\author{Alex Lin \quad (supervisor: Natalia Nolde)}
\date{August 18, 2026 \quad (version 8)}

\begin{document}
\maketitle


\section{Data and the two conditioning structures}\label{sec:data}

We have twelve hydrometric stations in British Columbia. For station $i$, $X_{i,t}$ is the daily streamflow on day $t$ and $B_{i,t}$ is the daily baseflow --- the slow, storage-fed component of the flow. Baseflow is measured directly, not derived by filtering the discharge series.

Three daily series exist per variable: the observed record 1968--2012 (about $45 \times 365.25 \approx 16{,}400$ days per station, complete at all twelve stations) and a climate-model series under each of two emission scenarios, RCP 4.5 and RCP 8.5, running the full 1968--2099 with a simulated historical segment of its own. 

The design decision is one fitting axis, 1968--2099, carrying the RCP 4.5 series alone --- the history record is not considered. The reason is a jump at the join: a simulated run evolves from its own initial conditions and is not calibrated to reproduce the observed history year by year, so concatenating the two sources produces an offset at 2012 and 2013, visible in baseflow, less obvious in streamflow.

RCP 8.5 enters later, as the second arm of the scenario comparison of Section~\ref{sec:rl}. The no-extrapolation logic is unchanged: the spline machinery of Section~\ref{sec:curves} cannot extrapolate --- a fitted curve is defined only where its basis has data --- so the forbidden move is continuing any curve past its data, and the axis must contain every date for which we ever want an answer. The future part of the axis is informed by scenario data, not by statistical extrapolation: extrapolating a fitted linear trend from the historical era alone is technically possible, but the observational era carries no information about how things evolve beyond it, and the uncertainty of such an extrapolation would be too large to use.

The covariate is a single running index $t$: day one of 1968 through the last day of the axis. Day-of-year and calendar year reappear only inside the return-level calculation, as the grouping of days into years. The seasonal cycle as a modelled shape is explained together with the basis in Section~\ref{sec:curves}, where the choice can be defended rather than just adopted.

Whatever else the conditional extremes model does, it estimates one target object:
\[
\Pr\bigl(X \le x \,\big|\, C = c\bigr) \quad \text{for } c \text{ large},
\]
where $C$ is the conditioning observable. The pairing convention is the definition of $C$: different conventions give different targets. Our project runs two structures.

Structure A, spatial conditioning. The random vector is streamflow at the twelve stations. After marginal transformation (Section~\ref{sec:common}), condition on station $i$ being large and model the rest:
\begin{equation}\label{eq:strA}
\mathbf{Y}_{-i} \,\big|\, Y_i = y \;\approx\; \mathbf{a}\, y + y^{\mathbf{b}}\, \mathbf{Z}, \qquad y > v,
\end{equation}
read componentwise. The scientific question is co-flooding: given an extreme at station $i$, what happens across the network. Note this is really twelve models, one per choice of conditioning station, fitted separately as in Heffernan and Tawn (2004), with self-consistency across them not automatic. Six pairs are already designed for reporting, from the dependence screening: (1) Barriere at the mouth given Barriere below Sprague; (2) Fraser at Hope given Fraser near Marguerite; (3) Hansard given McGregor; (4) Bridge given North Thompson; (5) Eagle given Illecillewaet; (6) Elk at Fernie given Smoky at Watino. 

The working order is one pair first: the full pipeline runs end to end on a single pair, and will repeat to see more results later.

Structure B, baseflow conditioning. The random vector is the pair $(B, X)$ at a single station: condition on $B$ large, model $X$. This is the $d = 2$ case of the identical machinery --- one conditioning direction, one scalar residual $Z$. The role mapping onto Jonathan et al. (2014) is exact: baseflow $B$ is the conditioning variate (their storm peak significant wave height), streamflow $X$ is the associated variate (their storm peak period), and time $t$ is the covariate (their storm direction). The scientific question is a compound-event one: given the catchment is in an extreme wetness state (baseflow is high), what is the distribution of streamflow.

The pairing convention for Structure B is decided: C1, fixed lag. The modelled vector is
\begin{equation}\label{eq:strB}
V_t^{(\ell)} = \bigl(B_{t-\ell},\; X_t\bigr), \qquad \ell \ge 1,
\end{equation}
with target $\Pr\bigl(X_t \le x \mid B_{t-\ell} = b,\ t\bigr)$ for large $b$: risk given the wetness state $\ell$ days earlier. The same-day convention was rejected because during a flood the same-day baseflow can reflect the event in progress, so ``given the catchment was already wet'' is not what it conditions on. Every dependence quantity now carries the lag as an index --- $a_\ell(t)$, $b_\ell(t)$, $G_\ell$ --- because each $\ell$ defines a different joint tail. {\color{blue}The conditioning value is the single lagged day $B_{t-\ell}$, as written; replacing it by a pre-event window average $\bar B_t^{(\ell,w)} = \frac{1}{w}\sum_{u=t-\ell-w+1}^{t-\ell} B_u$ --- which would change little for slow-varying baseflow but would protect against isolated measurement noise --- is an option still to be discussed, kept on the open list, not part of the current convention.} Since $a_\ell$ is estimable on a grid $\ell \in \{0, 1, \ldots, L\}$, plotting $\hat a_\ell$ against $\ell$ is the natural diagnostic for how fast the predictive content of wetness decays --- and if $\ell$ is picked from this plot, that selection step belongs inside the bootstrap loop, or the uncertainty bands will be too narrow. The event-level alternative (the transplant of Jonathan et al.'s storm-peak construction, declustering built into the data definition) falls with the declustering route, which Section~\ref{sec:rl} sets aside.

\section{Marginal model, one variable at a time}\label{sec:marg}

For a given station and a given date, the marginal model answers ``how unusual is this value for this station on this date.'' The peaks-over-threshold means do not model the whole distribution, model only its tail, because only the tail obeys a universal approximation. Under non-stationarity a fixed threshold is inappropriate as there is a high quantile in the dry season but a mediocre one in the wet season, so its exceedances are mostly ``wet-season events,'' not ``extreme events.'' The repair is to fix the probability and let the level move.

\subsection{Moving threshold}\label{sec:threshold}

The threshold is the conditional quantile curve $u(t; \tau)$ with $\Pr\bigl(X_t \le u(t; \tau)\bigr) = \tau$ at every date, so that ``exceedance'' means large relative to its own date. It is fitted by minimizing the penalized quantile-regression criterion
\begin{equation}\label{eq:qr}
\ell^{*}_{QR} = \Bigl\{\, \tau \sum_{i:\, r_i > 0} |r_i| \;+\; (1 - \tau) \sum_{i:\, r_i < 0} |r_i| \Bigr\} + \lambda_u R_u, \qquad r_i = x_i - u(t_i; \tau),
\end{equation}
where the bracketed part is the pinball loss (the tilted absolute-error criterion whose minimizer is the $\tau$-quantile) and $\lambda_u R_u$ is a roughness penalty (Section~\ref{sec:curves}). This is done simultaneously for an increasing series $\tau_1 < \cdots < \tau_D$ under the non-crossing constraint $u(t; \tau_d) \le u(t; \tau_{d+1})$ for every $t$, solved as one linear program. The working level $\tau^{*}$ (provisionally 0.95) is chosen by the covariate version of threshold stability: fit the GP above each rung, watch where its parameter estimates stabilize, take the lowest adequate rung.

\subsection{Generalized Pareto above the threshold}\label{sec:gp}

Above its own date's threshold, the exceedance is generalized Pareto with covariate-dependent parameters:
\begin{equation}\label{eq:gp}
\Pr\bigl(X_t > x \,\big|\, X_t > u(t; \tau^{*}),\ t\bigr) = \biggl(1 + \xi(t)\, \frac{x - u(t; \tau^{*})}{\sigma(t)}\biggr)_{\!+}^{-1/\xi(t)}.
\end{equation}
Read the two functions physically. The scale $\sigma(t)$ is the typical size of an excess at date $t$; a trend in it says flood overshoots are growing. The shape $\xi(t)$ is the heaviness of the tail at date $t$, and letting it move is a much stronger claim --- not ``floods got bigger'' but ``the character of the flood tail changed.'' The menu across the source papers: Jonathan et al. (2014) let $\xi$ vary smoothly; Faulkner et al. (2019) hold it constant, citing estimation noise; Lee and Lee (2019) go further and set $\xi = 0$. The decision is constant $\xi$ per station. Fitting is by maximum likelihood --- each exceedance contributes a GP log-density at its own covariate value --- plus roughness penalties on the parameter functions; positive parameters are built on the log scale so positivity is automatic (the geometric paper's style; a presentation-level difference from Jonathan et al., who impose $\sigma > 0$ directly). How every parameter curve here is represented --- constant, line or spline --- and the three decisions attached to each are collected in Section~\ref{sec:curves}, because the same machinery serves the dependence curves of Section~\ref{sec:dep}.

\section{Common scale: transform and checks}\label{sec:common}

The two fitted pieces assemble into a complete distribution function for each variable at each date --- the quantile ladder below the threshold, the GP tail above:
\begin{equation}\label{eq:F}
F(x; t) =
\begin{cases}
\tau_{d-1} + (\tau_d - \tau_{d-1})\, \dfrac{x - u(t; \tau_{d-1})}{u(t; \tau_d) - u(t; \tau_{d-1})}, & u(t; \tau_{d-1}) \le x < u(t; \tau_d),\\[12pt]
1 - (1 - \tau^{*}) \biggl(1 + \xi(t)\, \dfrac{x - u(t; \tau^{*})}{\sigma(t)}\biggr)_{\!+}^{-1/\xi(t)}, & x > u(t; \tau^{*}).
\end{cases}
\end{equation}
The below-threshold line is their linear interpolation between the fitted quantile curves. {\color{blue} In the stationary 2004 paper the below-threshold body is the empirical distribution function: sort the sample, and $F$ becomes a staircase that jumps by $1/n$ at each observation and is flat in between --- exact at the data points, discrete everywhere, and blind to $t$ by construction, since it pools all days into one common distribution. Jonathan et al.'s below-threshold piece (their Section 3.3, ``in the absence of a parametric form for the cumulative distribution function, we approximate it'') is, in essence, that empirical function with the steps smoothed: when $x$ falls between two grid levels, the probability is not read off a step but spread linearly along the slope between the levels, with the GP tail taking over above the threshold. The rungs are the fitted curves $u(t; \tau_d)$, the step heights are the fixed grid probabilities $\tau_d - \tau_{d-1}$, and at each fixed $t$ the interpolation fills linearly within each rung interval --- the ladder doing double duty, top rung as threshold, whole stack as the body of the distribution, replacing the pooled empirical function that is no longer legitimate because it ignores $t$.} 

Each observation is then transformed by the fitted probability integral transform {\color{green}to the standard Laplace scale (Keef et al., 2013, their equation 2.1)},
\begin{equation}\label{eq:pit}
{\color{green}
y_t =
\begin{cases}
\log\bigl(2\, F(x_t; t)\bigr), & F(x_t; t) < \tfrac{1}{2},\\[4pt]
-\log\bigl(2\,\{1 - F(x_t; t)\}\bigr), & F(x_t; t) \ge \tfrac{1}{2}:
\end{cases}}
\end{equation}
its probability under its own day's distribution, converted to {\color{green}standard Laplace units --- the distribution with density $\tfrac{1}{2} e^{-|y|}$, symmetric about zero with an exactly exponential tail on each side}. For every $t$, $y_t$ is a draw from one fixed distribution. {\color{green}The Laplace scale is the working scale because of those two tails: one parametric form in Section~\ref{sec:dep} then covers positive and negative association at once, where the Gumbel scale of the 2004 paper --- exponential upper tail, much lighter lower tail --- represents positive association only. The Gumbel version keeps a place in the report as the 2004 paper's own convention and as a sensitivity refit; every operational step in this document is on the Laplace scale.}

The first diagnostic comes before any slicing and before the {\color{green}Laplace} step: the fitted values $F(x_t; t)$ themselves, collected into a histogram, should look like draws from $\mathrm{Unif}(0,1)$ --- the flat histogram is the same statement as the {\color{green}Laplace} QQ line one transformation earlier, and a systematic bump or hole in it says the marginal model has missed something. This is the gate of stage 1: the transform and everything downstream wait until the marginal diagnostics pass and are documented.

The transformed sample is assumed to be marginally stationary. The check must be sliced. Otherwise, if the true distributions are higher in June and lower in February but the fitted $F(\cdot\,; t)$ sits in the middle of both, June values come out systematically too large on the transformed scale and February values too small --- and pooled together the two errors cancel, so the pooled sample looks {\color{green}Laplace} while the promise is false. 

A slice is a group of days sharing one calendar label, formed using exactly the labels the model chose to smooth over: season slices test the no-seasonal-term decision (Section~\ref{sec:seasonknots}), decade slices test whether the slow curves tracked the drift. For each slice, sort its transformed values and plot the $k$-th smallest against {\color{green}the standard Laplace quantile of $k/(n_S + 1)$ --- with $q = k/(n_S + 1)$, that is $\log(2q)$ when $q < \tfrac{1}{2}$ and $-\log\bigl(2(1 - q)\bigr)$ when $q \ge \tfrac{1}{2}$ --- a QQ plot against the standard Laplace}; points on the 45-degree line mean the slice passes. The numeric companion is to fit the two {\color{green}Laplace} parameters per slice by maximum likelihood and check they sit near $(0, 1)$ (their slices were consecutive windows on the {\color{green}same} Laplace scale; season and decade as the groupings is our adaptation). Sample sizes are comfortable: the axis has about $132 \times 365.25 \approx 48{,}200$ days, thirteen decade slices of roughly $3{,}700$ days each and season slices of roughly $12{,}000$ days; at $\tau^{*} = 0.95$, about $2{,}400$ exceedance days per station, roughly $180$ per decade.

The tail check is simple because of one fact: the threshold is the fitted $\tau^{*}$-quantile curve, so $F(u(t; \tau^{*}); t) = \tau^{*}$ on every day, and the moving threshold becomes one flat number on the {\color{green}Laplace scale, $y^{*} = -\log\bigl(2(1 - \tau^{*})\bigr)$}, the same for every day of every year. So ``above the threshold'' means the same thing in every slice, and two checks follow. Counts: within each decade slice the fraction of days above $y^{*}$ should be close to $1 - \tau^{*}$ --- this tests whether the threshold curve tracked the trend, and it is precisely the property Lee and Lee state as the purpose of a moving threshold. Within season slices the fractions will not be equal and are not supposed to be: on a snowmelt river the exceedances concentrate in the freshet, the accepted consequence of defining extreme by absolute discharge --- report the seasonal counts as description, not failure. Shapes: a QQ plot of only the exceedances in each slice against the {\color{green}Laplace tail above $y^{*}$ --- exactly unit exponential for $y^{*} \ge 0$, so the excesses $y - y^{*}$ are compared to a standard exponential}; winter slices may have almost no points, which is itself information. One refinement connecting to stage 2 (our reasoning, flagged): what the dependence model actually consumes is the days on which the conditioning variable is extreme, so the decisive version of the check is the sliced QQ restricted to those days --- if the transform is right where stage 2 lives, the no-seasonal-term decision holds for the dependence model even if the deep body of February is imperfect.

Two threshold levels exist and should not be confused. The marginal level $\tau^{*}$ lives on the original scale and its curve must move with $t$. The dependence threshold lives on the {\color{green}Laplace} scale, where all dates are already identically distributed, so a single constant $v$ (equivalently one non-exceedance probability $\kappa$, {\color{green}$v = -\log\bigl(2(1 - \kappa)\bigr)$ for $\kappa > \tfrac{1}{2}$}) is correct --- mapped back to the original scale it is secretly a moving curve. Both levels are selected by stability ladders. {\color{green}A third level rides along with the dependence stage and no ladder selects it: the constraint level $v_K$ at which the Keef et al.\ (2013) parameter constraints of Section~\ref{sec:dep} are imposed --- a single value on the Laplace scale, set above the largest observed value of the conditioning variable, because the constraints police extrapolation beyond the data; imposing them at the dependence threshold itself, $v_K = v$, gave poor fits in their examples (their Section 3.2).}

\section{Parameter curves: one formula, three decisions}\label{sec:curves}

Every parameter in this document is a curve in $t$: the threshold, scale and shape of the marginal stage (Section~\ref{sec:marg}), and later the dependence and residual curves of Section~\ref{sec:dep}. The machinery is common to all of them: first the one representation every curve shares, then what that representation is allowed to see, then the three decisions attached to every curve.

\subsection{One formula for every parameter curve: constant, line, spline}\label{sec:basis}

The functional-form decision is not made once for the whole model; it is made once per parameter curve --- $u$, $\sigma$, $\xi$, and later $a$, $b$ and the residual location and scale. Every candidate is the same formula with a different basis:
\begin{equation}\label{eq:basis}
\eta(t) = \sum_{k=1}^{p} \gamma_k\, B_k(t),
\end{equation}
where the $B_k$ are fixed building-block functions and the $\gamma_k$ are the only things estimated. One column ($B = [\,1\,]$) is the constant --- the stationary model. Two columns ($B = [\,1 \;\; t\,]$) is the straight line --- Faulkner et al., Lee and Lee. Cubic B-spline bumps at knots spread evenly over the axis give the flexible curve --- Jonathan et al., with periodicity switched off. The bases are nested: the constant is a special case of the line, the line of the spline.

{\color{blue} 
Axis A, the running date — the one formula  \eqref{eq:basis} actually uses. One point per calendar date, from the start of 1968 to the end of 2099. Every point occurs exactly once, so nothing can be fitted pointwise.
Axis B, day-of-year — the hypothetical one. The covariate value of an observation is only which day of the year it fell on; the year is thrown away. This axis has only 366 points.  For instance, every year of the record contains a June 14, but on this axis June 14 of one year and June 14 of another are different covariate values, so they are not repeats of each other. A day-of-year covariate would create exactly the missing repeats, by merging the years: every year's June 14 would sit at the single covariate value "June 14", and each calendar day's distribution could then be estimated directly from its own observations, with no further assumption. 

Axis B design is rejected, because merging the years deletes the year: June 14 early in the record and June 14 late in it would be forced to share one distribution, and change over the record — the object of study — could not appear. The running axis is chosen deliberately because nothing repeats on it; the years still supply the observations — tens of thousands of days — but some structure must convert them into a statement about one specific date. That structure is the smoothness assumption, and the basis  \eqref{eq:basis} is its concrete form.


The straight line $\eta(t) = \gamma_1 + \gamma_2 t$ is formula \eqref{eq:basis} with two columns. Every day contributes a fitting term containing $\gamma_1$ and $\gamma_2$, both numbers are estimated from the whole sample, and the value at any date is read off the fitted line — the date supplied one observation, the two numbers producing its value were estimated from everything, and nothing was ever fitted pointwise. Two features of the line are worth naming, because the spline keeps the first and changes the second: it is one single curve defined over the whole axis, and its influence pattern is global — every day moves the curve everywhere, since both coefficients sit in every day's fitting term. 

The spline keeps the mechanism and the one-curve property, and changes exactly the influence pattern: it makes influence local. A knot is a time point chosen before fitting — a design choice, not a parameter — and a set of knots placed along the axis cuts it into consecutive intervals. A basis function $B_k(t)$ is then a fixed, fully known function of $t$, built from cubic polynomial pieces joined smoothly at the knots: exactly zero over most of the axis, nonzero only over a stretch of a few consecutive knot intervals (four, for cubic splines), where it rises smoothly to one peak and returns to zero — the bump.

The bump's shape and position depend only on the knot locations, never on the data; one bump sits at each position along the axis, each shifted over by one knot, so neighbouring bumps overlap. Two consequences carry the whole construction: at any single date only a few bumps are nonzero (exactly four in the cubic case), and under any single bump lie all the days of the several years it spans.

The curve is the weighted sum \eqref{eq:basis}: at a given date most terms vanish, so $\eta(t)$ is a combination of the four covering bumps — a weighted average of their coefficients, in fact, since the covering bump values at any date sum to one — and each coefficient $\gamma_k$ is the local height of the curve over its bump's stretch of years: raising it lifts the curve there and changes nothing far away, while the overlaps join the stretches into one smooth curve over the whole axis, exactly as for the line. The data never move or reshape the bumps; they only choose the heights.

Estimation is penalized maximum likelihood. Each day's term contains $\eta$ at that day, hence only the few covering coefficients, so $\gamma_k$ appears in exactly the fitting terms of the days under bump $k$ — thousands of days at the threshold stage, which uses every day, still hundreds of exceedance days at the GP stage — while days outside its stretch produce terms in which $\gamma_k$ never appears and so cannot move its estimate (the roughness penalty links neighbouring coefficients, so a little influence extends one bump over; nothing reaches across decades). 
The opening problem — one observation per date, so nothing can be fitted pointwise — is now resolved. The day-of-year axis would have resolved it with repeats: many observations at one covariate value. The running axis resolves it with neighbours: many observations at nearby covariate values. Neighbours can replace repeats only because of the smoothness assumption: since the curve cannot change quickly, observations at nearby dates carry information about the date between them. The two models now compare cleanly. In both, the whole curve is estimated from all the years. They differ at a single date: the line produces its value there from two numbers that every day of the record helped estimate, while the spline produces it from four coefficients that only the days of the nearby years helped estimate. The regression form shows the resulting sample sizes directly. The fit is a regression whose design matrix has one column per bump; the entry in row $i$ and column $k$ is $B_k(t_i)$, the value of bump $k$ on day $i$. Each column has thousands of nonzero rows, one for every day under its bump, and there are only a handful of columns in total. A handful of unknowns, each estimated from thousands of days.}

The stationary count makes the budget concrete: two GP parameters per margin, the two dependence parameters $a$ and $b$, and the residual's location and scale --- eight numbers, the threshold not counted as a parameter --- against the roughly $48{,}200$ days of the axis. In the spline version each number becomes a curve and the count becomes the basis dimension per curve, still small against the sample; the count is an order-of-magnitude feasibility argument, not a complexity budget.

\subsection{What the curves can see: seasonality outside, knots as resolution}\label{sec:seasonknots}

There is a property of the single-$t$ choice which is crucial, and it is the cost of the count just given: the coefficient count stays small only because the knots sit years apart, and knots years apart also fix what the fitted curves are able to represent. The knot spacing sets the finest time pattern the basis can produce. Bumps that each span several years give a curve that can follow the slow drift across decades, but inside any single year such a curve can only drift as well: it cannot rise and come back down within twelve months, so it assigns June and February of the same year essentially the same parameter values, even though the flows in those months differ systematically in every year of the record. Every fitted $\eta(t)$ is therefore a trend curve, and seasonality is not merely absent from the fit but unrepresentable by the basis. Making the single-$t$ basis represent it would require knots every few weeks along the whole axis — thousands of coefficients, each estimated only from the few months of days under its own bump, and at the GP stage only from the few exceedances among those days: the feasibility argument of the previous paragraph, failing. And the expense would buy the cycle in the most wasteful way available: the seasonal cycle has essentially one shape, repeated every year, and a single-$t$ basis has no way to impose that repetition, so it would re-learn the same shape separately in each of the 132 years; the design that learns one cycle from all years at once is a periodic day-of-year basis — the merging of the years rejected above as a replacement for the trend axis, where it deletes the year, but suited to a component that repeats by definition. The model does neither and deliberately leaves seasonality outside the parameter curves. What makes the exclusion tolerable is the target: the annual return level is a whole-year summary, determined by how many exceedances the year produces and how large they are; where within the year they fall does not enter the annual answer — the argument is given in full in Section 6 — and yearly counts and sizes are exactly what a trend-only model tracks. The exclusion is a decision, not an oversight, and it is checked: the seasonal diagnostic of Section 3 tests whether it has damaged the fit, and if it fails, the repair is the design just named — swap the single-$t$ basis inside formula (7) for one that also contains the shared day-of-year cycle alongside the trend, with coefficients, penalty, regression form and fitting all unchanged, which is what makes the repair non-structural.

{\color{blue}
The paragraph breaks are mine; merge them into one block if the section wants it. The formula reference and "Decision 2" cross-reference carry over unchanged.
The template is Jonathan et al.'s own arrangement, verified against the paper: cubic B-splines at $p$ uniformly spaced knot locations (their Section 3.5); a roughness penalty built on first differences of neighbouring coefficients, $k=1$ throughout their work, in the penalized-spline tradition of Eilers and Marx (2010); and cross-validation estimating one number only, the roughness weight $\lambda$ — ten-fold in their simulation case, where the basis has $p = 60$ knots. The knots themselves are never selected by the data: count and placement are fixed before fitting, and the only data-driven smoothing choice is $\lambda$. Within one fit, several parameter curves share that single $\lambda$: each curve's roughness has its own coefficient in principle, but the relative sizes are fixed in advance (their Section 3.4, to reduce computational burden, the ratios set by experimentation), so one number is estimated per fit — the same single-$\lambda$ practice Decision 2 below assumes. (One nuance, kept for accuracy: the paper states this explicitly for the quantile-regression and conditional-extremes stages; for the GP stage it lists two roughness coefficients without saying they are combined.)

Uniform placement is usually a compromise, because data can be dense in some stretches of the covariate and sparse in others, leaving some coefficients with few observations under their bumps. On our axis it is exact rather than a compromise: the axis carries one observation per date, so equally spaced knots put equally many days between consecutive knots — spacing uniform in $t$ and spacing uniform in observation count are the same thing here. One step further, our inference, flagged: at the GP stage only exceedance days carry information, and the moving threshold — the covariate-dependent quantile at level $\tau^{*}$ — is built to hold the exceedance probability near $1-\tau^{*}$ on every date, so even the exceedance days should fall close to evenly along the axis. This is also why Jonathan et al.'s one refinement is dropped rather than copied. Storm directions bunch heavily in some sectors, so before placing uniform knots they move to a rank-transformed covariate, $\theta_i^{*} = \tfrac{360}{n}\,(r(\theta_i)-1)$ — that is, each value replaced by its position in the sorted list, spread evenly by construction (their Section 5.1). With our density already uniform, the rank transform reduces to the identity: the refinement does nothing here, and its cost never arises.

The working numbers, defaults until the representation itself is confirmed: spacing of about five years, which across the 132-year axis gives on the order of thirty coefficients per curve — the "basis dimension per curve" of the feasibility count above. Quantities that are selected from the data must be re-selected inside every bootstrap resample, precisely because they are data-selected — Jonathan et al.'s own bootstrap repeats the full analysis, including the cross-validation of $\lambda$, in each resample, and the lag, the smoothing weights and the model sizes join it here for the same reason. 
A sensitivity check is used for knots: refit with the spacing halved and with it doubled, confirm the fitted curves barely move — once the penalty is doing the smoothing, it is $\lambda$, not the knot count, that controls the effective flexibility, provided the basis is rich enough — and report the outcome in one sentence. }

\subsection{Three decisions per curve}\label{sec:threedec}

Three separate decisions attach to every parameter curve. 

Decision 1, model size: which family in \eqref{eq:basis} the curve may use --- one column, two columns, or the spline basis. Decision 2, smoothness: if the spline family is used, how strongly its roughness is penalized. Decision 3, software: which package computes the fit. The discussion below runs Decision 2, then Decision 3, then Decision 1: move two of the model-size procedure reads the penalized fit that Decision 2 defines, and move one leans on the bootstrap machinery that Decision 3 supplies, so the two supporting decisions come first. 

\subsection{Decision 2: smoothness}\label{sec:smooth}

Smoothness is enforced by adding $\lambda R(\gamma)$ to every fitting criterion, with $R$ built from differences of neighbouring coefficients. Two mechanisms exist for $\lambda$, interchangeable here and only here: cross-validation --- Jonathan et al.'s stated procedure, our choice for now --- or automatic selection as in the geometric paper, qgam's calibrated Bayesian tuning for the quantile stage and evgam's REML for the tail stage. {\color{blue}One verified software fact belongs here: no package performs this cross-validation for our criteria, so the loop is project code, and because of the serial dependence of Section~\ref{sec:rl} its folds must be contiguous blocks of days rather than random days --- random folds leak information across the fold boundary through the dependence; the qgam and evgam mechanisms are the built-in alternative precisely because they replace the loop, not because they implement it.}

{\color{green}What the qgam mechanism is and what it costs, since it is the leading candidate to replace the loop (Fasiolo et al., 2021a for the method, 2021b for the package). The pinball loss is replaced by a smooth approximation (their Extended Log-F loss), close enough to a likelihood that likelihood machinery applies: the roughness weight $\lambda$ is then selected by an approximate marginal likelihood, the standard automatic criterion in additive models, and the smoothness of the loss itself is selected by a separate calibration step, tuned so that the resulting intervals for the quantile curve hold roughly their advertised coverage. Three costs are recorded, and the first is load-bearing. First, each quantile level is fitted independently, with no non-crossing constraint --- the package paper itself states that fitted quantiles often cross --- so qgam can replace the loop for the single working curve $u(t; \tau^{*})$, while a qgam-fitted ladder loses exactly the non-crossing property the linear program of Section~\ref{sec:threshold} builds in, and would need a separate check or repair. Second, the calibration theory assumes independent observations; under the serial dependence of Section~\ref{sec:rl} the advertised coverage is a working approximation, so the block bootstrap stays the uncertainty engine whichever mechanism selects $\lambda$. Third, their own warning at extreme levels: when $\min\{n\tau,\ n(1-\tau)\}$ is small against the number of model coefficients, the fit leans on an estimated density and its bias needs checking --- comfortable at $\tau^{*} = 0.95$ with about $2{,}400$ exceedance days per station, a real cost only for the highest rungs of a ladder.}

Two limits keep this decision from swallowing Decision 1. First, every mechanism above answers ``what predicts best,'' and no prediction criterion produces a significance statement, so no choice of $\lambda$ can certify that a trend exists. Second, the penalty has a null space --- the shapes assigned zero roughness, never shrunk however large $\lambda$ grows: constants under a first-difference penalty, straight lines under a second-difference one --- so the smoothing machinery can flatten bends out of a curve but cannot test whether a null-space component belongs in the model at all. Both jobs fall to Decision 1.

\subsection{Decision 3: software}\label{sec:software}

For the threshold, three routes compute the same object: quantreg (the linear fit, with bootstrap standard errors), and qgam or evgam's asymmetric-Laplace (ald) family (the smooth fits, with automatic $\lambda$). For the GP stage: evgam, or hand-coded penalized likelihood. For the dependence stage: hand-coded, since no package among the source papers covers it. Routes are never selected by cross-validation or AIC --- that vocabulary belongs to the other two decisions. The rule is agreement: run the threshold routes side by side on the same task, check the fitted curves agree, then standardize on one, chosen by what the surrounding steps need --- quantreg supplies the tested linear fit that move one below requires; qgam and evgam supply the smooth fits that move two reads.

{\color{blue}The software claims in this document are verified against the packages' own sources, with the details and the full no-package list in the companion audit file; four findings matter enough to sit here. For the non-crossing ladder, two engines exist beyond hand-coding: quantregGrowth's gcrq fits penalized spline quantile curves for a whole grid of levels with non-crossing built in, and quantreg's constrained method imposes exact non-crossing inequalities within the linear family --- so the simultaneous fit of Section~\ref{sec:threshold} has native support, and a hand-coded linear program is a choice rather than a necessity. For the dependence stage, the hand-coded fitter carries a built-in unit test: run with all four curves forced constant, it must reproduce texmex, the reference implementation of the stationary 2004 model, whose own simulator (mexMonteCarlo) is the stationary version of the Section~\ref{sec:rl} recipe. texmex also defaults to Laplace margins with {\color{green}the Keef et al.\ (2013) constraints --- the working scale of Section~\ref{sec:common} and the constraints of Section~\ref{sec:dep} --- so the constant-curve unit test runs against the reference implementation on its own default settings}. And for the serial-dependence treatment, evgam ships a family that fits the extremal index inside the same additive framework --- a covariate-dependent $\theta(t)$ if the clustering itself drifts --- alongside the stationary estimators in exdex and extRemes.}

\subsection{Decision 1: model size --- the fixed two-move procedure, per parameter}\label{sec:size}

Move one, constant versus linear: fit both and test whether the trend coefficient earns its place --- likelihood-ratio test ($\Lambda = 2(\hat\ell_{\mathrm{lin}} - \hat\ell_{\mathrm{const}})$ against $\chi^2_1$) or AIC in the stages that have a likelihood (the GP margin; the dependence stage too, with the caveat that its likelihood is the working Gaussian one of Section~\ref{sec:dep}, so the test there is an approximate guide and the bootstrap interval for the trend coefficient is the honest backup) --- and, for the threshold, the slope judged against its bootstrap standard error from the quantreg machinery, since the pinball loss is not a likelihood and no likelihood-ratio test exists there: precisely Lee and Lee's device. Move one is the only step in the pipeline that can produce the sentence ``a trend exists and is significant'' --- the claim the paper's title rests on, and the reason Decision 1 cannot be delegated to any smoothing mechanism. 

{\color{blue}Move one answer the question of ``does a trend exist'' , which is the hypothesis $\gamma_2 = 0$ --- and which tool tests it is decided by the stage's fitting criterion, because tests are built on criteria. Where the criterion is a genuine likelihood, at the GP margin, fit both candidates, record each maximized log-likelihood $\hat\ell$ (the log of how probable the data are under the fitted model), and read $\Lambda$: if the true slope were zero, $\Lambda$ would behave like a draw from $\chi^2_1$, one degree of freedom because the models differ by one parameter, and a $\Lambda$ too large for that distribution says the slope earned its place --- AIC is the same comparison with a fixed charge per extra parameter in place of a significance level. Where the criterion is the working Gaussian likelihood, at the dependence stage, the same $\Lambda$ can be formed but only trusted as an approximate guide, and the honest question is whether the slope's block-bootstrap interval excludes zero. Where there is no likelihood at all, at the threshold stage --- the pinball loss is a loss function, not a probability model, so $\hat\ell$ does not exist and $\Lambda$ cannot be formed --- the identical question is answered by the slope against its bootstrap standard error. 

One disambiguation, because the pipeline contains {\color{green}four} objects the word threshold could point to, and move one tests exactly one of them. The threshold that enters move one is the moving quantile curve $u(t; \tau^{*})$ of Section 2.1: the station's fitted high-quantile curve through time. It is a parameter curve like $\sigma(t)$ or $\alpha(t)$, and it is tested like one. It is not the level $\tau^{*}$. The level is a probability — a single design number, chosen once by the stability ladder — and it has no time axis, so there is no trend in it to test. It is also not the dependence threshold $v$ of Section 3. That is a single constant on the {\color{green}Laplace} scale, where the marginal transformation has already removed the time variation from the margins, so there is no curve to test there either. {\color{green}And it is not the Keef constraint level $v_K$ of Sections 3 and 5: a single design value on the Laplace scale, placed above the largest observed conditioning value, again with no time axis and nothing to test.} For the threshold curve, move one therefore asks a concrete question: does this station's high quantile drift along the axis, or is one flat quantile enough for the whole record — the constant $u(t) = \gamma_1$ against the line $u(t) = \gamma_1 + \gamma_2 t$, with the slope judged against its bootstrap standard error, as stated above. The constant is allowed to win, and a win for the constant is a legitimate outcome, not a contradiction of this section's opening. The opening's claim is conditional: a fixed threshold is wrong when non-stationarity is present. Move one is the per-station check of exactly that condition. Where the constant wins, the check has found no detectable trend in the high quantile, so the condition is not met, and a flat threshold there is not the mistake the opening warns against. What the opening forbids is assuming stationarity without checking it; a flat threshold adopted because the check found no trend is a conclusion, not an assumption.
}

Move two, linear versus spline, needs no separate test, because the spline contains the line and the penalized fit performs the comparison itself: fit the penalized spline and read whether it collapses --- effective degrees of freedom near 1 and a visually straight curve mean the data have flattened the spline into the line, and the line is reported with the spline kept as a robustness figure. Faulkner et al.'s empirical finding is worth holding as an expectation: the formal tests tend to keep the simpler model; over the 132-year axis, though, a straight line is much less likely to survive against the spline than it would have over 45 observed years.

Two safeguards close the procedure. First, the post-fit reading: once everything is fitted, look at the curves themselves --- a parameter whose fitted spline comes out essentially flat falls back to a constant, whichever move produced it, because the flexible model is safe precisely in the sense that its worst outcome is a flat curve. Second, the stationary competitor: the all-constant model is fitted alongside the chosen one on the same data, and the diagnostics of Sections~\ref{sec:common} and~\ref{sec:dep} arbitrate between them --- non-stationarity is a claim to be won against a fitted baseline, not a default.

\section{Dependence model}\label{sec:dep}

Writing $(Y_C, Y_X)$ for the {\color{green}Laplace}-scale pair --- conditioning variate and associated variate, in either structure --- the non-stationary Heffernan--Tawn model is
\begin{equation}\label{eq:ht}
Y_X \,\big|\, \bigl(Y_C = y,\ t\bigr) \;=\; a(t)\, y \;+\; y^{b(t)} \bigl\{ m(t) + s(t)\, Z \bigr\}, \qquad y > v, \quad Z \sim G.
\end{equation}
Four smooth functions of the covariate are: $a(t) \in {\color{green}[-1, 1]}$ and $b(t) \in {\color{green}(-\infty, 1)}$ are the dependence parameters, and $m(t)$, $s(t)$ are the location and scale of the standardized residual. When the conditioning variable sits at rarity $y$, the associated variable typically matches the fraction $a(t)$ of that rarity, plus scatter of size growing like $y^{b(t)}$. The division of labour is worth one sentence: in the $y$ direction the normalizing functions $a(t)\, y$ and $y^{b(t)}$ are specified parametric forms --- assumptions plugged in, not objects of estimation --- while in the $t$ direction the four curves are exactly what estimation is about. {\color{green}The range $[-1, 1]$ for $a$ is a property of the Laplace working scale: both tails of the Laplace are exponential, so the single form \eqref{eq:ht} covers positive association ($0 < a \le 1$) and negative association ($-1 \le a < 0$) at once (Keef et al., 2013, their equation 2.3); the Gumbel convention of the 2004 paper restricts $a$ to $[0, 1]$ and represents positive association only.} In Structure A the equation is read componentwise as \eqref{eq:strA}, giving eleven curve sets $\bigl(a_{j|i}(t), b_{j|i}(t)\bigr)$ per conditioning station; in Structure B it is the $d = 2$ case with the lag index attached, $a_\ell(t), b_\ell(t)$.

The load-bearing assumption of the whole non-stationary extension, stated in their Section 3: the standardized residual has a limit distribution that does not depend on the covariate,
\[
\Pr\bigl(Z \le z \,\big|\, Y_C = y,\ t\bigr) \;\to\; G(z), \qquad y \to \infty,
\]
for a non-degenerate $G$. It is checkable on our data: after removing $a, b, m, s$, the standardized residual cloud should look homogeneous across decades --- if it does not, the covariate effect has not been fully captured by the four parameter functions. The model's main working diagnostic is the same object seen locally: plot $\hat Z$ against the conditioning value above $v$; any drift in level or spread says the normalization failed. Two further cautions carry over: the exceedance days feeding the fit are treated as independent, which multi-day flood clusters violate --- the serial-dependence treatment of Section~\ref{sec:rl} is the answer; and everything is per conditioning direction --- the pair for ``$j$ given 1 extreme'' is not the pair for ``1 given $j$ extreme,'' and models fitted in different directions are not guaranteed mutually consistent (the self-consistency issue from the 2004 discussion).

The dependence class is never assumed --- it is read off the fitted $\hat a(t)$. At $a = 1, b = 0$ the model says $Y_X \approx Y_C + Z$: asymptotic dependence, the link that never fades ($\chi > 0$){\color{green}; $a = -1, b = 0$ is its exact negative mirror}. Everything {\color{green}strictly between is asymptotic independence, the sign of $a$ giving the direction of association --- the Gaussian copula sits at $a = \operatorname{sign}(\rho)\, \rho^2, b = 1/2$}; exact independence is $a = 0, b = 0$. Because $a$ is a function of the covariate, the class is even allowed to differ across covariate values.

Case 1, Barriere given Barriere --- shared floodwater --- holds $\chi(0.95)$ between 0.94 and 0.99 in every six-year window, so its fitted $\hat a$ must land near 1, or the fit is wrong before the data are. Cases 2, 3 and 5 fill the middle territory ($\chi(0.95)$ roughly 0.7--0.87, 0.67--0.82, 0.53--0.77): pull that is real but sub-proportional, where $0 < a < 1$ earns its keep --- and notably Case 5 crosses river networks entirely, so the divider between corners is not network membership but how much meteorology and water two gauges share. Case 6 is held back pending a data-version discrepancy. Case 4, Bridge given North Thompson, is the pair that breaks the stationary picture: statistically indistinguishable from independence in 1947--52, roughly $\chi \approx 0.4$ with the interval clear of zero by 1995--2006 --- read through the Gaussian benchmark, an effective $a$ drifting from about 0 toward about 0.2 (a gloss, flagged). The one pair at the boundary is the only pair that moves, and it is the concrete, data-given argument for letting $(a, b)$ vary with the covariate. The empty fourth corner is its own finding: nothing in this network is independent of anything else.


Structure and estimation. In Structure A, $\mathbf{Y}_{-1} \mid Y_1$ is one joint model whose location-and-spread part decomposes into the eleven pair fits, and whose jointness is carried entirely by the eleven-dimensional residual vector $\bigl(Z_{2|1}, \ldots, Z_{12|1}\bigr)$ with one unknown joint distribution $G_{|1}$: on a day when the same storm pushed several stations beyond their fitted lines at once, that shared overshoot lives nowhere in the $a, b$ curves --- only in the joint law of the residuals.

Estimation uses a working Gaussian likelihood for $Z$, penalized by roughness terms with the weights chosen as in Section~\ref{sec:smooth}, fitted pair by pair; the fitted residuals are then kept, one row per historical extreme day, as an empirical sample from $G$.

{\color{green}Two Keef et al.\ (2013) facts attach to this fit. First, the scoring device: they tried an independent-Laplace working likelihood against the independent-Normal one and found the Normal better across a broad range of simulated examples (their Section 2.4), so the change of margins does not change the working likelihood --- it stays Gaussian. Second, the maximization is constrained. Writing, for the days above $v$, $z(q)$ for the empirical $q$-quantile of the model's own fitted residual, and $z^{+}(q)$, $z^{-}(q)$ for the empirical $q$-quantiles of $Y_X - Y_C$ and $Y_X + Y_C$ --- the residuals the data would show under exact asymptotic positive dependence ($Y_X = Y_C + Z^{+}$) and its exact negative mirror ($Y_X = -Y_C + Z^{-}$) --- their ordering constraint (their equation 3.1), stated here in the stationary parameters, is
\begin{equation}\label{eq:keef}
-\,y + z^{-}(q) \;\le\; \alpha\, y + y^{\beta}\, z(q) \;\le\; y + z^{+}(q) \qquad \text{for all } y > v_K \text{ and all } q \in [0, 1].
\end{equation}
Read plainly: beyond the data, the fitted conditional quantiles may not sit above what perfect positive dependence would give, nor below what perfect negative dependence would give --- the model is forbidden to claim a link stronger than perfect in either direction, which is exactly the kind of inconsistent extrapolation an unconstrained fit can produce. Their Theorem 1 converts the ordering into explicit feasibility conditions on $(\alpha, \beta)$ at the level $v_K$; checking the two endpoint levels $q = 0$ and $q = 1$ was found to cover all $q$ between (their Section 3.2); and the implementation is blunt: the working likelihood is set to zero wherever the conditions fail, so the maximizer can never land outside the feasible set. The level $v_K$ sits above the largest observed conditioning value, as Section~\ref{sec:common} records. Carrying the ordering into the non-stationary fit --- imposed at each $t$ through the curves $a(t), b(t)$ and the per-$t$ residual quantiles --- is our extension of their stationary constraint, flagged.}

{\color{blue}Spelled out, because this table is the model's memory. Take one historical extreme day --- a real calendar day on which the conditioning variable sat above $v$, at {\color{green}Laplace} value $y$. For each associated station $j$, equation \eqref{eq:ht} is solved backwards for that day: subtract the fitted level $a_{j|1}(t)\, y$, divide by the fitted growth factor $y^{b_{j|1}(t)}$, remove the residual location $m_{j|1}(t)$ and divide by the residual scale $s_{j|1}(t)$; the number left over is that station's standardized residual for that day --- how far above or below its fitted expectation the station sat, in standard units, with the effects of rarity $y$ and covariate $t$ both removed. Eleven such numbers, laid side by side, are one row. The rows stack over every historical extreme day into a table with one row per extreme day and one column per associated station, and that table, kept intact, is the estimate of $G$: the working Gaussian was only ever a scoring device --- the residual is treated as normal to score fits against each other, never believed --- so $G$ carries no formula and is represented by the rows themselves, each row one observed draw from the joint law. The licence for reusing a 1975 row at any other $t$, including a future one, is exactly the load-bearing assumption above: after the four curves remove the covariate effect, $G$ does not depend on $t$ --- which is what the sliced residual check verifies. In Structure B a row is a single number, and row-resampling reduces to the ordinary draw of residuals with replacement; the row language earns its keep when there are several columns.} At simulation time whole rows are resampled --- the same day's pattern for all stations at once, never independent draws from separate columns --- because the rows are the only place the model stores the fact that stations overshoot together; independent draws would silently destroy every joint-risk output while leaving all single-pair outputs looking unchanged. {\color{blue}One caution rides on the rows and lands in Section~\ref{sec:rl}: a six-day event above $v$ contributes six near-duplicate rows, overweighting one storm in the fit --- another face of the serial dependence whose price is paid by the block bootstrap.} The deliverable at the end of this stage is exactly 2004 Figure 6 translated into hydrology: conditional return values of streamflow, given baseflow at a level with non-exceedance probability 0.99 or 0.999, as a function of $t$; the generative recipe that produces it is written out in Section~\ref{sec:rl}.

\section{Return level and the 2099 axis}\label{sec:rl}

The return level is where the calendar reappears --- and only as the grouping of days into years. Nothing new is fitted. Under non-stationarity there is no single ``$T$-year level'' per station; the definition must be indexed by a reference year (the convention from the Faulkner discussion). Fix a station, a period $T$ and a reference year $Y$: freeze the fitted curves as they stand in year $Y$ (essentially constant within any single year, since every curve is a slow trend curve), let $M_Y$ be the annual maximum of a year with those parameters, and define $x_T(Y)$ by
\begin{equation}\label{eq:rl}
\Pr\bigl(M_Y > x_T(Y)\bigr) = \frac{1}{T}.
\end{equation}
Read plainly: a year with the climate of year $Y$ exceeds this level at least once with probability $1/T$. Under a trend only this probability reading survives; the waiting-time reading (``we will wait about $T$ years'') is no longer true, because real future years will not behave like year $Y$. Every return-level statement in the paper is a probability statement about a reference year, not a forecast --- one sentence in the Results section, then the figures speak.

Construction from the daily model. For $x$ above the threshold on day $t$, the daily exceedance probability is the chance of being an exceedance day at all, times the GP tail beyond $x$,
\begin{equation}\label{eq:pxt}
p(x, t) = (1 - \tau^{*}) \biggl(1 + \xi\, \frac{x - u(t)}{\sigma(t)}\biggr)_{\!+}^{-1/\xi},
\end{equation}
and the year chains its days together, treating days as independent for the moment:
\begin{equation}\label{eq:annual}
\Pr\bigl(M_Y \le x\bigr) = \prod_{t\, \in\, \text{year } Y} \bigl(1 - p(x, t)\bigr).
\end{equation}
This is also where the trend-only frame is internally consistent: the annual answer depends on how many exceedances a year produces and how large they are, which the trend-only model tracks, while the within-year timing --- what a seasonal curve would have described --- cancels out of the annual product. The caution that does not cancel is serial dependence: consecutive flood days are not independent draws. The treatment is decided, and it is the modern one --- keep every day and adjust the inference, rather than decluster first (merging consecutive exceedance days into one event and keeping its peak, the older route): the extremal-index correction (Lee and Lee) enters the chaining here, and every uncertainty statement comes from the block bootstrap, which resamples whole blocks of consecutive days so that the dependence inside a block is preserved rather than destroyed.

{\color{blue}The confusion worth preventing here is between two questions that sound alike. Whether each day has its own marginal distribution is one question --- answered yes by the whole of Sections~\ref{sec:marg} and~\ref{sec:common}, every date owning its $F(\cdot\,; t)$ --- and whether knowing today's value changes what to expect tomorrow is a different question entirely, a property of days jointly; independence is the second question, and separate marginals say nothing about it. June 3 and June 4 of a freshet each carry their own, slightly different distribution, and knowing June 3 ran enormous still makes an enormous June 4 far more likely. The data settle the second question by themselves: rivers have storage, snowmelt and storms last days, today's water feeds tomorrow's gauge, so exceedances arrive in runs --- a six-day freshet above threshold is six exceedance days from one physical event --- and dependence between consecutive days is an empirical fact that no modelling choice can switch off. The model treats days as independent in exactly two formulas, as a working assumption: the fitting criteria, where each exceedance day contributes its own multiplied term, and the chaining product \eqref{eq:annual}. Writing the assumption down does not make the data obey it, the same way the working Gaussian never made the residual normal, and each of the two uses is paid for where it bites. The chaining is paid by the extremal index. If the year really were 365 fresh chances, \eqref{eq:annual} would be exact; with clustering the year holds fewer effective chances than days, and the extremal index $\theta \in (0, 1]$ measures exactly this --- $1/\theta$ is the average run length among extreme days. Lee and Lee's own statement of the classical result, verified against the paper: with a common marginal $F$, the maximum of the independent series has limit $G(x)$ exactly when the maximum of the dependent series has limit $G^{\theta}(x)$; ``the extremal index $\theta$ is equal to 1 for a perfectly independent series, and less than 1 for a dependent series,'' introduced because return-period estimates ``could be severely biased if the data are dependent.'' Carried into the per-day construction, the corrected chaining is
\begin{equation}\label{eq:theta}
\Pr\bigl(M_Y \le x\bigr) \;\approx\; \Bigl[\prod_{t\, \in\, \text{year } Y} \bigl(1 - p(x, t)\bigr)\Bigr]^{\theta}.
\end{equation}
At $\theta = 1$ nothing changes; at $\theta = 1/3$ extremes come in three-day runs on average and the year holds about a third as many independent opportunities. The direction of the error if $\theta$ is ignored follows from the exponent: the independent product makes ``exceed somewhere this year'' look easier than it is, overstating the annual exceedance probability and inflating the return levels --- wrong even when the error lands on the cautious side. The inference is paid by the blocks: dependent days carry less information than independent ones, so standard errors, the move-one tests and likelihood intervals computed under independence come out too narrow --- too confident --- and the block bootstrap restores honesty by resampling whole blocks of consecutive days and refitting everything per replicate, the dependence inside a block preserved in every resample. That pairing, extremal index plus bootstrap, is Lee and Lee's own pipeline, and it is the sense in which keeping all days is the opposite of declaring them independent: declustering manufactured approximate independence by collapsing each run to its peak, at the price of discarded data and an arbitrary event boundary, while the decided route keeps the dependent days, admits the dependence, and attaches the two corrections exactly where the independence assumption is used.}

Computation, two routes with the same answer for one station. Route A, direct: substitute the fitted curves into \eqref{eq:annual} and solve \eqref{eq:rl} for $x$ by one-dimensional root finding, over stations, years and the chosen $T$. Route B, simulation: generate many synthetic years under the fitted model for reference year $Y$ and read off the empirical $(1 - 1/T)$ quantile of the simulated annual maxima. The name matters: generating from the fitted model is Monte Carlo --- approximating quantiles and probabilities by large-scale sampling from a fitted distribution --- not the bootstrap, which resamples observed data and is reserved here for uncertainty; the residual draw inside the generation resamples fitted residuals with replacement, and that is the only bootstrap-like step. {\color{blue}A third term stays separate from both: a simulation study --- data generated from curves whose truth is chosen and known, the whole pipeline run on them, the estimates compared to the planted truth, as Jonathan et al. do before touching their sea data --- is a validation exercise, distinct from this generation step, where nothing is known and everything is estimated; and the climate-model series is simulated data as input, a third thing again.} Route B is the workhorse --- it is Jonathan et al.'s own route for return values, and the only route once a question involves more than one variable, because then the answer has no formula. Uncertainty comes from the block bootstrap: refit everything on block-resampled data --- including the selection steps, the lag $\ell$, the smoothing weights, the model sizes --- and take the spread of the recomputed curves.

The figure: one panel per station, reference year on the horizontal axis, flow in real units on the vertical, one curve per return period with bootstrap bands, plus the companion table of start-of-axis versus end-of-axis levels with intervals. Vertical reading: fix a year, read the flood size at a given rarity. Horizontal reading: fix a flow level and watch which curves the horizontal line crosses as the years advance --- the flood that was a 100-year event early in the record becoming, say, a 50-year event later is the single most quotable sentence the figure can produce. A rising curve is intensification; a flat curve is also a finding (a stationary model would have sufficed there, and this framework is what makes that checkable); a falling curve is physically possible on snowmelt rivers.

The single-station curve is a stage-1 product: it depends only on that station's own marginal model, and all twelve panels could be computed before any dependence is fitted. Stage 2 earns its keep the moment the question crosses variables: given station A exceeds its own $T$-year level for year $Y$, the distribution of station B; the probability that both exceed their levels in the same year, and whether it has grown; and the Structure-B deliverable, conditional return values of flow given extreme antecedent wetness. All are read off the Route B simulation with row-at-a-time residual sampling.

The stage-2 quantities have their own definition, and it is per day: each date carries its own joint distribution --- 2027 is one distribution, 2100 another --- so the conditional return level is indexed by $t$ exactly as the marginal level is indexed by its reference year. For the Structure-B pair,
\begin{equation}\label{eq:condq}
Q_t\bigl(\tau \,\big|\, B_{t-\ell} \ge b\bigr) \;=\; F^{-1}_{\,X_t \,\mid\, B_{t-\ell} \ge b}(\tau),
\end{equation}
the $\tau$-quantile of streamflow on day $t$ given that baseflow $\ell$ days earlier exceeded the level $b$: the return-level idea carried past one dimension, with the conditioning event written into the symbol. Conditioning on $B_{t-\ell} = b$ in place of $\ge b$ works the same way, and the unconditional marginal quantile of $X_t$ is the one-dimensional companion, reportable beside it. No closed form exists for \eqref{eq:condq}, so the estimate is generative, resting on the stochastic representation \eqref{eq:ht}: draw the conditioning variable from its fitted marginal restricted to the exceedance region, draw the residual with replacement from the kept rows, assemble the associated variable through $a$, $b$, $m$, $s$ at the chosen $t$, repeat on the order of ten thousand times, and take empirical quantiles of the generated conditional sample --- with enough repetitions that some hundreds to a thousand of the generated points land above the target quantile, the order-of-magnitude rule that makes an empirical high quantile usable. The same generated cloud delivers any functional: the conditional mean, the conditional quantile of \eqref{eq:condq}, and the probability of any region of interest, read off as the fraction of generated points falling inside. The 2004 paper's own simulation figure is this method at work --- the conditional cloud comes from the fitted model extrapolating, not from the data --- which is why the Section~\ref{sec:dep} deliverable is available at every $t$ on the axis.

Finally, the RCP 4.5 axis makes the scenario stage literal: the same fitted curves reach 2099 because the basis has data there, so $x_T(Y)$ and the conditional quantiles are produced for future dates with no extrapolation ever performed --- with the honest note that every diagnostic now tests agreement with the scenario series, observations having stayed outside the fit. The analysis then runs on two axes: evolution along $t$ within a scenario, and differences between scenarios, with the whole pipeline refitted on RCP 8.5 and the two answers set side by side --- ``how much worse'' is the main line, and the start-versus-end companion table is its tabular form. A stationary baseline completes the comparison: either the stationary model fitted to the historical segment, where no real trend lives, or the full model with $a$ and $b$ held constant, the latter carrying clear model bias. For display, fix one year --- 2050, say; any year works --- and plot the estimates across it: choosing a date to display is illustration, not selection bias, provided the questions the figure answers are written down first. And none of it is a literal prediction: the whole exercise is scenario analysis --- statements about how flood risk behaves under a specified emissions pathway, read comparatively.

\section{Open decisions (to discuss)}\label{sec:open}

Every item on this list carries a default --- either a value already set or a decided procedure whose outcome arrives mechanically once the code runs --- and the defaults are sufficient: on them alone the whole analysis runs end to end and the report can be finished, so nothing below blocks completion. Open to discussion therefore does not mean undecided. It means the default is the normal setting, and this is the short list of places where a credible alternative exists and switching to it could change the results --- which is exactly why the items stay visible: a change that moves the answers is a finding about sensitivity, worth a robustness check or a paragraph in the discussion, and some alternatives are worth mentioning in the report even when never adopted {\color{green}(the Gumbel working scale is the clearest case: nothing operational runs on it, yet it keeps a place in the report as the 2004 paper's own convention and as a sensitivity refit beside the Laplace fit)}. Each item below states its default beside what remains open.

\begin{enumerate}
\item Threshold levels: marginal $\tau^{*}$ (provisionally 0.95) and dependence $\kappa$, both by stability ladders. Default: $\tau^{*} = 0.95$; for $\kappa$ no provisional number is set, so its default is the procedure itself --- the lowest adequate rung the ladder returns.
\item {\color{blue}Structure B timing: the lag $\ell$ (grid up to $L$, the $\hat a_\ell$ decay plot, selection inside the bootstrap); and whether to replace the single lagged day by a pre-event window average at all --- to be discussed --- and, if adopted, the window $w$.} Default: the single lagged day under convention C1 of Section~\ref{sec:data}, with $\ell$ read off the $\hat a_\ell$ decay plot inside the bootstrap; no window average.
\item Per-parameter model sizes (Decision 1): the two-move procedure is fixed; the outcomes (which parameters end constant, linear, spline) await the code. Default: constant $\xi$ per station is already fixed (Section~\ref{sec:gp}); every other curve takes whatever size the two moves return.
\item Smoothing mechanism (Decision 2): cross-validation now; qgam / evgam (REML) as the automatic alternative{\color{green}, the qgam mechanism and its three costs stated in Section~\ref{sec:smooth}}. Software (Decision 3): standardize on one threshold route after the quantreg / evgam-ald / qgam agreement check. Default: cross-validation with contiguous-block folds (Section~\ref{sec:smooth}); for software, run the agreement check and standardize on whichever route the surrounding steps need.
\item Sharing across stations: common forms or common parameters versus twelve free copies (notes page pending). Default: twelve free copies --- the per-station fitting the pipeline already runs.
\item The moving-window screening: the window width for the windowed $\chi$ and $\bar\chi$ sequences. Default: the six-year width the screening of Section~\ref{sec:dep} already used.
\end{enumerate}

\section*{Main references}

Heffernan, J. E. and Tawn, J. A. (2004). A conditional approach for multivariate extreme values. \textit{JRSS B} 66, 497--546.

{\color{green}Keef, C., Papastathopoulos, I. and Tawn, J. A. (2013). Estimation of the conditional distribution of a multivariate variable given that one of its components is large: additional constraints for the Heffernan and Tawn model. \textit{J. Multivariate Analysis} 115, 396--404 (Laplace working scale; the unified form $a(y) = \alpha y$, $b(y) = y^{\beta}$ with $\alpha \in [-1, 1]$; the parameter constraints; the Normal-over-Laplace working-likelihood finding).}

Jonathan, P., Ewans, K. and Randell, D. (2014). Non-stationary conditional extremes of northern North Sea storm characteristics. \textit{Environmetrics} 25, 172--188.

Faulkner, D., Warren, S., Spencer, P. and Sharkey, P. (2019). Can we still predict the future from the past? Implementing non-stationary flood frequency analysis in the UK. \textit{J. Flood Risk Management} (model-size comparison; per-reference-year reporting).

Lee, J. and Lee, T. (2019). Extreme snow events in New York City (linear-quantile threshold with bootstrap standard errors; extremal index; $\xi = 0$ variant).

{\color{green}Fasiolo, M., Wood, S. N., Zaffran, M., Nedellec, R. and Goude, Y. (2021a). Fast calibrated additive quantile regression. \textit{JASA} 116, 1402--1412 (the qgam method: smoothed pinball loss, automatic smoothing weight, calibrated intervals).}

{\color{green}Fasiolo, M., Wood, S. N., Zaffran, M., Nedellec, R. and Goude, Y. (2021b). qgam: Bayesian nonparametric quantile regression modeling in R. \textit{J. Statistical Software} 100(9) (the package; source of the quantile-crossing caveat).}

Murphy-Barltrop, C., Wadsworth, J., de Carvalho, M. and Youngman, B. (the geometric-approach paper on non-stationary extremal dependence; per-window stationarity check; qgam / evgam smoothing machinery; log-scale convention for positive parameters).

\end{document}